\documentclass[11pt]{article}
\usepackage[margin=1in]{geometry}
\usepackage{amsmath,amssymb,amsthm,booktabs,array,longtable}
\usepackage[T1]{fontenc}
\usepackage[utf8]{inputenc}
\usepackage{lmodern}
\usepackage[colorlinks=true,urlcolor=blue,citecolor=blue,linkcolor=blue]{hyperref}
\setlength{\emergencystretch}{1em}
\newtheorem{theorem}{Theorem}
\newtheorem{corollary}[theorem]{Corollary}
\newcommand{\PP}{\mathcal P}
\newcommand{\NN}{\mathcal N}
\newcommand{\G}{\mathcal G}
\newcommand{\even}{\mathrm{even}}
\newcommand{\odd}{\mathrm{odd}}
\title{Ba\v{s}i\'c--Gottlieb--Krnc (2026), Section 5, Question 1 (p.~17):\\an alternative four-part proof after Mehat (2026),\\with explicit strategies and small-row stacking}
\author{Haobo Ma\\\small ChronoAI Pte Ltd, Singapore\and
Wenlin Zhang\\\small National University of Singapore\\\small ChronoAI Pte Ltd, Singapore}
\date{6 November 2026}
\begin{document}
\maketitle
\begin{abstract}
For Ba\v{s}i\'c--Gottlieb--Krnc (2026), Section~5, Question~1, p.~17,
we give an alternative proof of Mehat's classification (23 August 2026)
of all losing four-part positions in normal-play CRIM, with no bound
on the part sizes.
For $a\ge b\ge c\ge d\ge1$, the largest part $a$ must be even;
four disjoint cases involving $d$, $c$ and the parity of $b$ give the remaining conditions.
The proof identifies every possible column move by its height and gives
an explicit winning move in each complementary case.
Consequences include even-rank necessity and hook-stacking closure with at most
four rows or columns, the odd-gap winning-move conclusion through six rows
or columns using Mehat's rank theorem, exact padding formulas for one- and two-part inputs,
and an explanation of 21 of the 38 diagrams in their Appendix~A.
\end{abstract}

\section{The classification}
In CRIM, a position is an integer partition
$\lambda=[\lambda_1,\ldots,\lambda_r]$ with
$\lambda_1\ge\cdots\ge\lambda_r>0$.
A move deletes one row or one column of its Young diagram and joins the
remaining pieces. Equivalently, a row move deletes a part, and a column move
conjugates the partition, deletes a part, and conjugates back.
The player with no move loses. Write $\PP$ for the losing positions and
$\NN$ for the winning positions. The empty partition $[\,]$ is in $\PP$.
Section~5, Question~1, p.~17 of Ba\v{s}i\'c--Gottlieb--Krnc~\cite{BGK}
asks: ``Which partitions consisting of four parts are losing?''
Mehat~\cite[Theorem~4, p.~2; proof in Section~8, p.~9]{Mehat}
answered this question in the manuscript included in his public archive
of 23 August 2026. His Theorems~1--2 classify positions through six rows
or columns by a finite table of truncated gap types. The following
equivalent four-case formula has a direct cell-count induction proof
and an explicit winning-move table.

\begin{theorem}[Mehat (2026), Theorem 4]\label{thm:main}
Let $a\ge b\ge c\ge d\ge1$. Then $[a,b,c,d]\in\PP$ if and only if
$a$ is even and the applicable row of the following table holds.
\begin{center}
\begin{tabular}{ll}
\toprule
Range & Additional condition\\
\midrule
$d\ge2$ & $b,c$ are not both odd\\
$d=1,\ c=1$ & $b=1$ or $b$ is even\\
$d=1,\ c=2$ & $b\ge3$ or $a=2$\\
$d=1,\ c\ge3$ & $b$ is odd\\
\bottomrule
\end{tabular}
\end{center}
\end{theorem}

For example, $[2,2,2,1]\in\PP$ but $[4,2,2,1]\in\NN$.
Also $[6,4,3,2]\in\PP$ whereas $[4,3,3,2]\in\NN$.
The boundary $[2,2,2,1]$ explains why an even largest part alone does not
resolve the cases ending in $[2,1]$.
For the previously studied almost-hooks $[a,b,b,b]$ with $b\ge3$,
the table gives losing exactly when $a,b$ are even, in agreement with
Ba\v{s}i\'c--Gottlieb--Krnc (2022), Theorem~5.1~\cite{BGK22}, restated
as Theorem~8 of~\cite{BGK}. The classification extends beyond that family.

\section{Proof and winning moves}
For a position with $r$ parts, put $\lambda_{r+1}=0$.
Let $R_i$ denote deletion of row $i$.
A column of height $t$ exists precisely when $\lambda_t>\lambda_{t+1}$.
Every such column has the same follower, denoted $C_t$: subtract one from
the first $t$ parts and discard zeros.
There are $\lambda_t-\lambda_{t+1}$ columns of height $t$, so the
$R_i$ and these $C_t$ exhaust all moves, including interior-column deletions.
The strict gap at $t$ preserves the order of the parts.
Every move decreases the number of cells and preserves the class of
partitions with at most four parts.

Define a predicate $P^*$ on that class as follows:
\begin{center}
\begin{tabular}{ll}
\toprule
Position & Condition for $P^*$\\
\midrule
$[\,]$ & always\\
$[a]$ & never\\
$[a,b]$ & $a$ is even\\
$[a,b,c]$ & $a>1$, $a,b$ odd, and $(c>1\text{ or }b=1)$\\
$[a,b,c,d]$ & the conditions of Theorem~\ref{thm:main}\\
\bottomrule
\end{tabular}
\end{center}
All displayed parts are positive and nonincreasing.
The two-part row is Proposition~10 of~\cite{BGK}.
The three-part row is the criterion of its Theorem~11 with the necessary
exception $[1,1,1]$: deletion of its only column reaches $[\,]$, so
$[1,1,1]\in\NN$. This precise correction was already given by
Mehat~\cite[Corollary~25, p.~9; Section~12, p.~12]{Mehat}.
We reprove the corrected three-part row together with
the four-part row, without relying on the omitted boundary case.

\subsection{No $P^*$ position has a $P^*$ follower}
The empty position has no follower.
From a two-part $P^*$ position, row deletion leaves one part.
Column deletion changes the even first part to an odd part and leaves one
or two positive parts; it cannot erase an even-width position completely.
Neither type of follower satisfies $P^*$.

From a three-part $P^*$ position, every row deletion leaves two parts
with odd first part. A column deletion retaining three parts changes the
first part to an even number. If the bottom row disappears, then $c=1$;
the predicate forces $b=1$, and only the positive row $[a-1]$ remains.
Again every follower fails $P^*$.

Now let $[a,b,c,d]$ satisfy the four-part predicate, so $a$ is even.
Every $R_i$ with $i>1$ leaves three parts with even first part.
For $R_1$ to satisfy $P^*$ we would need
\[
 b>1,\quad b,c\text{ odd},\quad (d>1\text{ or }c=1).
\]
If $d>1$ this contradicts the first row of Theorem~\ref{thm:main};
if $c=1$ it contradicts the second row.
Every legal $C_t$ with $t<4$ leaves four positive parts with odd first
part, as does $C_4$ if $d>1$.

It remains to consider $C_4$ when $d=1$.
Its follower is the positive part of $[a-1,b-1,c-1]$.
If $c=1$, it has one or two positive parts and odd first part.
If $c\ge2$ and $b$ is odd, it has three parts and even second part.
Otherwise the four-part predicate forces $c=2$ and either $b\ge3$ or
$a=2$. For $b\ge3$, the follower has last part $1$ and second part
at least $2$, so it is not a three-part $P^*$ position.
For $a=2$ we have $b=c=2$, and the follower is $[1,1,1]$,
also outside $P^*$.
This exhausts all possibilities.

\subsection{Every other position has a $P^*$ follower}
A one-part position deletes its row to $[\,]$.
For two parts outside $P^*$, $a$ is odd.
If $a=1$, use $C_2$ to reach $[\,]$.
If $a>1$ and $a>b$, use $C_1$; if $a=b>1$, use $C_2$.
The latter followers have two positive parts and even first part.

For three parts outside $P^*$, if $a$ is even use $R_2$;
if $a$ is odd and $b$ even use $R_1$.
Both followers have two parts and even first part.
Otherwise $a,b$ are odd.
If $a=1$, use $C_3$ to reach $[\,]$.
The only remaining failure of $P^*$ is $a>1$, $b>1$, $c=1$.
Use $C_3$ to reach the two-part position $[a-1,b-1]$ with even first part.

For four parts outside $P^*$ the following thirteen cases are disjoint
and exhaustive. Each row assumes $a\ge b\ge c\ge d\ge1$ and failure of
the four-part predicate.
\begin{center}
\small
\begin{tabular}{rlll}
\toprule
 & Conditions & Move & Follower\\
\midrule
1 & $a$ even, $d>1$ & $R_1$ & $[b,c,d]$\\
2 & $a$ even, $d=c=1$ & $R_1$ & $[b,1,1]$\\
3 & $a$ even, $d=1<c$ & $C_4$ & $[a-1,b-1,c-1]$\\
4 & $a=1$ & $C_4$ & $[\,]$\\
5 & $a>1$ odd, $b=c=1$ & $C_1$ & $[a-1,1,1,1]$\\
6 & $a$ odd, $c=1<b$ & $R_2$ & $[a,1,1]$\\
7 & $a,b$ odd, $c>1$ & $R_4$ & $[a,b,c]$\\
8 & $a,c$ odd, $b$ even, $d>1$ & $R_2$ & $[a,c,d]$\\
9 & $a$ odd, $b,c$ even, $d>1$ & $C_1$ & $[a-1,b,c,d]$\\
10 & $a,c$ odd, $b$ even, $d=1<c$ & $C_2$ & $[a-1,b-1,c,1]$\\
11 & $a$ odd, $b,c$ even, $d=1$, $c\ge4$ & $C_3$ & $[a-1,b-1,c-1,1]$\\
12 & $a$ odd, $b=c=2$, $d=1$ & $C_3$ & $[a-1,1,1,1]$\\
13 & $a$ odd, $b>2$ even, $c=2$, $d=1$ & $C_1$ & $[a-1,b,2,1]$\\
\bottomrule
\end{tabular}
\end{center}
Here is a check of every row, including the existence of the chosen column.
In row~1, failure of the predicate forces $b,c$ odd, while $d>1$,
so the follower satisfies the three-part rule.
In row~2 it forces $b>1$ odd, giving the three-part hook.
In row~3, $b$ is even. If $c\ge3$, the follower has odd first two parts
and last part at least $2$. If $c=2$, failure of the predicate forces
$b=2$ and $a>2$, giving $[a-1,1,1]$ with $a-1>1$ odd.
In all these cases $C_4$ exists because $d=1$.

Row~4 is the all-one column.
In row~5, $a>1=b$ ensures $C_1$ exists; its follower satisfies the
$c=d=1$, $b=1$ case.
Row~6 gives a three-part hook, and row~7 gives three parts with
odd first two entries and third entry greater than $1$.
Row~8 has the same three-part property.
In row~9, parity implies $a>b$, so $C_1$ exists and its follower
has even first part, $d>1$, and even middle parts.
In row~10, parity implies $b>c$, so $C_2$ exists;
the follower has even first part, odd second part, and $c>1$.
In row~11, $c>1=d$ ensures $C_3$ exists; the new second part
is odd and the new third part is at least $3$.
In row~12, $c=2>d$ ensures $C_3$ exists, and the follower has
even first part and all other parts equal to $1$.
Finally in row~13, $a>b$ follows from parity; the follower satisfies
the $c=2,d=1,b\ge3$ row of Theorem~\ref{thm:main}.

To check exhaustion, first separate even and odd $a$.
For even $a$, split $d>1$, $d=c=1$, and $d=1<c$.
For odd $a$, separate $a=1$, then $c=1$, then $c>1$ with $b$ odd
or even. In the even-$b$ branch, split $d>1$ from $d=1$ and then
split the parity and size of $c$ as in rows~8--13.
Thus every position outside $P^*$ has a legal $P^*$ follower.

\subsection{Induction}
Induct strongly on the number of cells.
Every follower has fewer cells, so assume its losing status agrees with
$P^*$. A $P^*$ position has no losing follower and is losing.
A position outside $P^*$ has a losing follower and is winning.
This proves the classification on the entire move-closed class of
partitions with at most four parts, and in particular proves
Theorem~\ref{thm:main}. In the standard mex description of finite
impartial games~\cite[\S6, Satz~I]{Sprague}, equivalently,
\[
 \G(\lambda)=\operatorname{mex}\{\G(\mu):\lambda\longrightarrow\mu\}
\]
is zero exactly when $P^*(\lambda)$ holds.

\section{Consequences and source coverage}
\begin{corollary}\label{cor:rank}
A losing CRIM position with at most four rows or at most four columns
has even rank $|r-\lambda_1|$.
\end{corollary}
\begin{proof}
For two rows the largest part is even; for three it is odd; for four
it is even. The empty position has rank zero and one-row positions
are winning. Conjugation preserves legal moves and exchanges the row
and column counts, so the conclusion also holds for at most four columns.
\end{proof}
This proves the corresponding restricted case of Conjecture~1
in~\cite{BGK}. Mehat~\cite[Corollary~3, p.~2; Section~9, p.~10]{Mehat}
proved the stronger statement with six in place of four; we use that
stronger result in the next deduction.
Conjugating Theorem~\ref{thm:main} also classifies all positions with
exactly four columns, by applying its table to their four column heights.

\begin{corollary}\label{cor:gap}
The conclusion of Observation~1 of~\cite{BGK} holds unconditionally
for positions with at most six rows or columns: from a winning position
of even rank, either the first row or the first column can be removed
to reach a losing position, and at least one of the first two row-length
or column-height gaps is odd (missing entries are zero).
\end{corollary}
\begin{proof}
The scope is closed under moves. A winning position has a losing follower;
Mehat's rank theorem makes that follower's rank even. Deleting a row
other than the first changes the row-count parity but not the width,
so cannot give this follower. Deleting the first row changes the width
to $\lambda_2$ and the row count to $r-1$; equality of the old and new
rank parities requires $\lambda_1-\lambda_2$ odd. A column deletion
that leaves the row count unchanged changes the width parity and cannot
give the follower either. To eliminate a row it must delete the first
column, with all bottom unit rows disappearing. The new row count is
$\lambda'_2$, so parity requires $r-\lambda'_2=\lambda'_1-\lambda'_2$
odd. These are all possibilities. The empty follower is included by
taking missing parts and column heights to be zero.
\end{proof}

\begin{corollary}\label{cor:pad}
Let $a\ge b\ge1$ and $\lambda=[a,b]$. For the padding operation
$\widehat\lambda=[a,a,a,b]$ of Section~5.1 of~\cite{BGK},
\[
 \widehat\lambda\in\PP
 \quad\Longleftrightarrow\quad
 a\text{ is even and }(b\ge2\text{ or }a=2).
\]
Moreover,
\[
 [a,b,1,1]\in\PP
 \quad\Longleftrightarrow\quad
 a\text{ is even and }(b=1\text{ or }b\text{ is even}).
\]
\end{corollary}
\begin{proof}
For $[a,a,a,b]$ with even $a$, if $b\ge2$ use the first row of
Theorem~\ref{thm:main}. If $b=1$ and $a=2$, use its third row;
if $b=1$ and $a\ge4$, its fourth row fails because its second part
is $a$, an even number. Odd $a$ always fails.
For $[a,b,1,1]$, use the second row of the theorem.
\end{proof}
Thus, among losing two-part positions, Question~5 has the exact answer
$b\ge2$ or $a=2$. For every two-part position the implication
$\widehat\lambda\in\PP\Rightarrow\lambda\in\PP$ in Question~6 holds.
The second formula answers the analogous two-part question about
appending two unit parts, and proves its converse implication on two-part inputs.

\begin{corollary}\label{cor:one}
For every $a\ge1$, both $[a,a,a]$ and $[a,1,1]$ are losing exactly when
$a>1$ is odd, whereas $[a]$ is always winning.
Consequently, among one-part inputs, the implication in Question~6 holds
exactly when $a=1$ or $a$ is even. The converse implication for appending
two unit parts fails for every odd $a\ge3$.
\end{corollary}
\begin{proof}
Substitute into the three-part row of $P^*$; delete the row from $[a]$.
\end{proof}

\begin{corollary}\label{cor:stack}
Conjecture~8 of~\cite[Section~5.1, p.~19]{BGK} holds whenever the resulting
stack has at most four rows or at most four columns: stacking a losing
partition inside a losing hook gives a losing partition.
\end{corollary}
\begin{proof}
Write the outer hook as $[c,1^{r-1}]$. A nonempty losing hook has $r,c\ge2$
of the same parity, by the source's Proposition~7. A stack inside it has
the form $[c,1+\mu_1,\ldots,1+\mu_k,1^{r-k-1}]$, where
$k\le r-1$ and $\mu_1+1\le c$. The empty inner partition leaves the
losing hook unchanged; a nonempty losing inner partition has $k\ge2$.
For $r\le4$ the possibilities are therefore $(r,k)=(3,2),(4,2),(4,3)$.
When $(r,k)=(3,2)$, $c$ is odd and $\mu_1$ is even, so the three-part
criterion holds. When $(r,k)=(4,2)$, $c$ is even and $1+\mu_1$ is odd;
the four-part criterion with last part $1$ holds. When $(r,k)=(4,3)$,
$c$ is even and $\mu_1$ is odd, so the second part is even and the last
part is at least $2$; again the criterion holds. Transposition exchanges
rows and columns and conjugates a stack to the corresponding stack of
the conjugate hook and inner partition, giving the column case.
\end{proof}
The row-removal argument in Observation~2, p.~20 of~\cite{BGK}, now
applies without a conjectural premise when its stack has at most four
rows or columns: the indicated successive even hooks give losing stacks
by Corollary~\ref{cor:stack}. Thus a stack of odd hooks ending in $[1]$
in this range is winning, and removal of its first row is a winning move.

\begin{corollary}\label{cor:boundary}
The formula in Conjecture~2 of~\cite[Section~5, p.~18]{BGK} fails at
$(r,k)=(1,0)$: the square rectair $R^0_{1,1}=[1]$ has Grundy value $1$,
not $2$. The corrected value at this boundary is $1$.
\end{corollary}
\begin{proof}
The sole distinct follower of $[1]$ is the empty partition, of value $0$,
so its mex is $1$. In the printed formula, neither $r$ even nor $k<r-1$
holds, and $r\notin\{3,5\}$, giving $2$.
\end{proof}
This corrects the boundary value, not the formula for $r\ge2$.
The winning status of $[1]$, and that of $[1,1,1]$ used above, also
follow from the source's Proposition~7 and the hook formula of
\cite[Theorem~4.1]{BGK22}; neither individual value is new.
The staircase conjecture in~\cite[Section~6, Conjecture~1]{BGK22}
also predicts $2$ instead of $1$ at $S_1=[1]$ and $S_3=[3,2,1]$.
Indeed $S_3$ has distinct followers $[2,1]$, $[3,1]$, $[3,2]$,
$[2,1,1]$ and $[2,2,1]$, of values $0,3,3,3,3$, respectively,
so its mex is $1$.
These exceptions are already included in the revised Conjecture~4
of~\cite{BGK}; its unrestricted odd-staircase formula is not proved here.

\paragraph{Appendix A, p.~24.}
Number the source's 38 diagrams in reading order, left to right in each
row. The 21 with four rows are numbers $1$--$12$, $17$--$24$, and $38$.
They are precisely
\[
\begin{gathered}
[4,2,2,2],\ [6,2,2,2],\ [6,4,1,1],\ [6,3,2,2],\ [6,4,2,1],\\
[4,4,3,3],\ [6,4,2,2],\ [6,6,1,1],\ [8,2,2,2],\ [8,4,1,1],\\
[6,4,3,2],\ [6,5,3,1],\ [8,3,2,2],\ [8,4,2,1],\ [8,6,1,1],\\
[10,2,2,2],\ [6,4,3,3],\ [6,4,4,2],\ [8,4,2,2],\ [8,5,2,1],\ [10,4,1,1].
\end{gathered}
\]
Substitution in Theorem~\ref{thm:main} explains each losing status.
The remaining 17 diagrams have more than four rows and columns; the
accompanying exact evaluator confirms their individual losing values.

The following table records the scope relative to the numbered future-work
items in~\cite{BGK}. ``Not covered'' concerns the scope of this note.
\begingroup\small
\begin{longtable}{p{0.29\textwidth}p{0.62\textwidth}}
\toprule
Source item & Coverage\\
\midrule
Question 1, Section 5, p.~17 & Answered: exactly Theorem~\ref{thm:main}.
Reproved; prior proof Mehat (2026), Theorem 4.\\
Conjecture 1, p.~17 & Reproved through four rows or columns; prior proof
Mehat (2026), Corollary 3, through six.\\
Observation 1, p.~17 & Proved through six rows or columns using Mehat's
Corollary 3: Corollary~\ref{cor:gap}.\\
Question 2 & Zero Grundy values determined for at most four parts;
nonzero values not covered in general.\\
Conjecture 2, p.~18 & Refuted (explicit counterexample $r=1,k=0$);
boundary corrected to $1$. Formula for $r\ge2$ not covered.\\
Conjectures 3, 4, 5, p.~18 & General Grundy-value formulas not covered.\\
Question 3 & General rectair Grundy values not covered.\\
Conjecture 6; Question 4 & Conway-pair classification not covered.\\
Conjecture 7 & General thick-hook melds not covered.\\
Question 5, Section 5.1, p.~19 & Answered for two-part losing inputs:
$b\ge2$ or $a=2$. No one-part losing inputs exist.\\
Question 6, p.~19 & Answered: yes for all two-part inputs;
for one-part inputs, yes exactly when $a=1$ or $a$ even.\\
Unit-part appending question, p.~19 & Answered on one- and two-part inputs:
Corollaries~\ref{cor:pad} and~\ref{cor:one}.\\
Conjecture 8; Observation 2, pp.~19--20 & Proved for resulting stacks with
at most four rows or columns; general statements not covered.\\
Conjecture 9 & General parity self-conjugate criterion not covered.\\
Proposition 10 & Reproved; prior proof Ba\v{s}i\'c--Gottlieb--Krnc (2026).\\
Theorem 11, p.~9 & Refuted-and-corrected at $[1,1,1]$; corrected criterion
reproved. Prior correction/proof Mehat (2026), Corollary 25 and Section 12.\\
Appendix A, p.~24 & 21 of 38 losing diagrams explained by the theorem;
17 other individual positions confirmed by finite checks.\\
BGK (2022), Theorem 5.1 & Prior proof cited: the almost-hook subfamily
$[a,b,b,b]$, $b\ge3$.\\
BGK (2022), Conjecture 1 & Refuted at $S_1,S_3$; prior correction
BGK (2026), Conjecture 4 cited. General formula not covered.\\
\bottomrule
\end{longtable}
\endgroup

\section{Verification coverage}
The Lean declaration
\texttt{D5.S3.Combinatorics.Games.}\allowbreak
\texttt{CrimFourPartClassification.result}, in the
\href{https://github.com/the-omega-institute/trureturing/blob/e0da8fcc80226febc4e7ea0836927bc71d158292/D5/S3/Combinatorics/Games/CrimFourPartClassification.lean}{pinned classification file},
has statement \texttt{claim}: for every $1\le d\le c\le b\le a$,
$\G([a,b,c,d])=0$ if and only if $a$ is even and
\begin{align*}
 &(d\ge2\ \text{and $b,c$ not both odd})\\
 {}\text{or }&(d=c=1\ \text{and $(b\text{ even or }b=1)$})\\
 {}\text{or }&(d=1,\ c\ge2,\ b\text{ odd})\\
 {}\text{or }&(d=1,\ c=2,\ (b\ge3\text{ or }a=2)).
\end{align*}
Splitting $c=2$ from $c\ge3$ makes this exactly
Theorem~\ref{thm:main}. The
\href{https://github.com/the-omega-institute/trureturing/blob/e0da8fcc80226febc4e7ea0836927bc71d158292/D5/S0/Certificates/Games/CrimGrundyRefutation.lean}{pinned evaluator file}
defines \texttt{rowMoves} and \texttt{conjugate}.
Its \texttt{moves} function combines row and column deletion;
\texttt{grundy\_eq} gives the mex recurrence for \texttt{grundy}.
Strict cell descent is given by \texttt{move\_decreases}.
Its separate \texttt{result} negates the printed Conjecture~3 using
$R^5_{7,6}$; that seven-row statement is outside the present note's coverage.
The consequence statements above have written proofs in this note.
The declaration \texttt{result} in the
\href{https://github.com/the-omega-institute/trureturing/blob/e0da8fcc80226febc4e7ea0836927bc71d158292/D5/S0/Certificates/Games/CrimRectairSquareBoundaryRefutation.lean}{pinned square-rectair file}
proves the negation of the printed universal Conjecture~2, using exactly
the boundary in Corollary~\ref{cor:boundary}; it asserts no replacement
formula for $r\ge2$.

The accompanying \texttt{verify.py} independently constructs all source
moves by conjugation and computes their exact mex values.
For all $46\,376$ partitions with at most four parts and largest part
at most $30$, it checks the classification, conjugation, move completeness,
strict descent, both $P^*$ transition conditions and the consequences above.
There are $40\,920$ four-part inputs, of which $15\,761$ are losing;
each of the thirteen strategy cases occurs.
These finite checks complement the unrestricted induction proof.
The same evaluator checks the one-part formulas through $a=30$, all
fitting losing-hook stacks with at most four rows and outer width at most
$30$ and their conjugates, and the 38 diagrams of Appendix~A separately.
They also give
$\G([1,1,1])=1$, $\G([4,3,3,2])=4$ and $\G([6,4,3,2])=0$.

\begin{thebibliography}{4}
\bibitem{BGK}
I.~Ba\v{s}i\'c, E.~Gottlieb and M.~Krnc,
\emph{CRIM: A Natural Game on Integer Partitions},
\href{https://arxiv.org/abs/2606.16828v1}{arXiv:2606.16828v1}, 15 June 2026.
Sections~2.1, 2.3, 3.1, 5 and 5.1.
\bibitem{Mehat}
S.~S.~Mehat,
\emph{Losing positions of CRIM for partitions with at most six parts},
14-page manuscript in the supplementary artifact, version 1.0.0,
23 August 2026, \href{https://doi.org/10.5281/zenodo.22070744}{Zenodo 22070744}.
Theorems~1, 2, 4; Corollaries~3, 25; Sections~8, 9 and~12.
\bibitem{BGK22}
I.~Ba\v{s}i\'c, E.~Gottlieb and M.~Krnc,
\emph{Some observations on the Column-Row game},
Proceedings of the 8th Student Computing Research Symposium (SCORES'22),
Ljubljana, 6 October 2022, pp.~29--32.
\href{https://zalozba.fri.uni-lj.si/SCORES2022.pdf}{Proceedings},
Theorems~4.1 and~5.1; Section~6, Conjecture~1.
\bibitem{Sprague}
R.~Sprague,
\emph{\"Uber mathematische Kampfspiele},
Tohoku Mathematical Journal, First Series \textbf{41} (1935), 438--444.
\href{https://www.jstage.jst.go.jp/article/tmj1911/41/0/41_0_438/_article}{Publisher text};
\S6, Satz~I, pp.~439--440.
\end{thebibliography}
\bigskip
\noindent\small\textcopyright\ 2026 Haobo Ma and Wenlin Zhang.
Licensed under Creative Commons Attribution 4.0 International
(CC BY 4.0).
\end{document}
